\documentclass[exercice,a4]{thedocument}%
\usepackage{texmaths-tikz}%
\usepackage{pstricks}%
\usepackage{pst-bezier}%
\usepackage{pstricks-add}%
\usepackage{pst-eucl}%
\usepackage{pst-func}%
\usepackage{pst-tree}%
\usepackage{pst-circ}%
\usepackage{tikz}%
\usepackage{pgfplots}%
\usepackage{tkz-tab}%
\usepackage{tikzpagenodes}%
\usepackage{tabularray}%
\usepackage{tabularx}%
\usepackage{multirow}%
\usepackage{multicol}%
\usepackage{scratch3}%
\usepackage{enumitem}%
\usepackage[np]{numprint}%
\usepackage{pas-tableur}%
\usepackage{listings}%
\usepackage{graphicx}%
\usepackage{adjustbox}%
\usepackage{needspace}%
\usepackage{xcolor}%
\usepackage{eurosym}%
\usepackage{diagbox}%
\usepackage{colortbl}%
\usepackage{multido}%
\usepackage{icomma}%
\usepackage[french]{babel}%
\graphicspath{{../brevet/2026/}}%
\startdatakeys
    \source{26GENMATPO3}
    \BO{2026}
    \cycle{4}
    \competencesDuSocle{CA,RA,CO}
    \objectifsApprentissage{C4A1aa,C4A2ad,C4B2cb,C4D2aa,C4D3ac}
    \motsCles{étendue,tableur,pourcentage}
\enddatakeys
\definecolor{gris}{gray}{0.5}%
\providecommand{\wideparen}[1]{\overset{\frown}{#1}}%
\begin{document}%
\providecommand{\degres}{^\circ}%
Voici un tableau des températures moyennes annuelles à Paris, mesurées par une station météorologique entre 2015 et 2024. Les températures sont exprimées en degré Celsius ($^{\circ}$C).

\begin{center}
\begin{tblr}{hline{1,3}=1pt,vlines,colspec={|c|c|c|c|c|c|c|c|c|c|c|}}
	\hline
Année & 2015 & 2016 & 2017 & 2018 & 2019 & 2020 & 2021 & 2022 & 2023 & 2024 \\
	\hline
Température ($^{\circ}$C) & 13,47 & 12,87 & 13,39 & 13,91 & 13,66 & 14,28 & 12,86 & 14,30 & 14,27 & 13,50 \\
	\hline
\end{tblr}
\end{center}
\small{Source : \url{https://opendata.paris.fr/explore/dataset/temperature_moyenne_paris/table/?sort=annee}}
%\normalsize

\begin{enumerate}
\item
	\begin{enumerate}
		\item Calculer l'étendue de la série des températures moyennes annuelles de 2015 à 2024.
		\item Recopier sur la copie la formule qui doit être saisie dans la cellule L2 pour obtenir la température minimale de 2015 à 2024.
		
\begin{center}
\begin{tblr}{*{4}{|X[c]}|}
			\hline
			\texttt{=MIN(B1:K1)} & \texttt{=MIN(B2:K2)} & \texttt{=MIN(B2:L2)} & \texttt{=MIN(B1:L1)} \\
			\hline
\end{tblr}
\end{center}
	\end{enumerate}
\item La température moyenne annuelle à Paris en 1924 était de $10{,}73^{\circ}$C tandis que celle de 2024 était de $13{,}50^{\circ}$C. Un élève affirme qu'en passant de 10,73 à 13,50, l'augmentation est supérieure à $25\,\%$. A-t-il raison ? Argumenter votre réponse en expliquant la démarche.
\end{enumerate}

Pour la suite du problème, on estime que la température moyenne annuelle à Paris augmente d'environ $0{,}025^{\circ}$C chaque année.

\begin{enumerate}
	\setcounter{enumi}{2}
	\item En 2024, la température moyenne à Paris était de $13{,}50^{\circ}$C. En se basant sur l'estimation annoncée ci-dessus, vérifier que la température moyenne à Paris en 2034 est estimée à $13{,}75^{\circ}$C.
	\item \textbf{(Question algorithmique)}\\
	L'objectif du programme ci-dessous, réalisé avec le logiciel Scratch, est d'estimer la température moyenne à Paris pour une année choisie.

	\begin{center}
	\begin{scratch}[scale=0.9,num blocks]
		\blockinit{Quand \greenflag est cliqué}
		\blocksensing{demander \ovalnum{Choisir un nombre supérieur à 2024} et attendre}
		\blockvariable{mettre \selectmenu{Nombre d'années d'écart} à \ovaloperator{\ovalvariable{réponse} - \ovalnum{2024}}}
		\blockvariable{mettre \selectmenu{Température} à \ovaloperator{\ovalnum{13.5} + \ovalnum{0.025} * \ovalvariable{Nombre d'années d'écart}}}
		\blocklook{dire \ovaloperator{regrouper\ovalnum{ La température en degré Celsius sera de} et {\ovalvariable{Température}}}}
	\end{scratch}
	\end{center}

	Un utilisateur choisit 2044 comme nombre de départ. Quelle valeur est alors stockée dans la variable \texttt{Nombre d'années d'écart} ?
	\item 
	\begin{enumerate}
		\item Résoudre l'équation : $13{,}5 + 0{,}025x = 15$.
		\item En quelle année la température moyenne annuelle à Paris sera-t-elle estimée à $15^{\circ}$C ?
	\end{enumerate}
\end{enumerate}

\needspace{10\baselineskip}
\medskip
\showthesource
\end{document}
